\documentclass[12pt,a4paper,psamsfonts]{amsart}
\usepackage{amssymb,mathtools}
\usepackage{calc}
\setlength{\topmargin}{0cm}
\setlength{\oddsidemargin}{0cm}
\setlength{\evensidemargin}{0cm}
\setlength{\marginparwidth}{0cm}
\setlength{\marginparsep}{0cm}
\setlength{\textheight}{\paperheight - 2in - 35pt}
\setlength{\textwidth}{\paperwidth - 2in}
\setlength{\headheight}{12.5pt}
\setlength{\headsep}{25pt}
\setlength{\footskip}{30pt}
\renewcommand{\baselinestretch}{1.3}
\pagestyle{headings}
\raggedbottom
\setlength{\emergencystretch}{2em}
\usepackage[colorlinks=true,urlcolor=blue,linkcolor=blue,citecolor=blue]{hyperref}
\newtheorem*{theorem}{Theorem}
\theoremstyle{definition}
\newtheorem*{setup}{Setup}
\theoremstyle{plain}
\newtheorem{lemma}{Lemma}[section]
\newtheorem{proposition}[lemma]{Proposition}
\newtheorem{corollary}[lemma]{Corollary}
\newcommand{\Q}{\mathbb Q}
\newcommand{\C}{\mathbb C}
\newcommand{\OO}{\mathcal O}
\newcommand{\mult}{\operatorname{mult}}
\newcommand{\nef}{\operatorname{n}}

\usepackage{tikz-cd}
\newcommand{\R}{\mathbb R}
\newcommand{\Z}{\mathbb Z}
\newcommand{\PP}{\mathbb P}
\DeclareMathOperator{\Pic}{Pic}
\DeclareMathOperator{\Supp}{Supp}
\DeclareMathOperator{\rk}{rk}
\DeclareMathOperator{\ord}{ord}
\DeclareMathOperator{\vol}{vol}
\DeclareMathOperator{\Sym}{Sym}
\DeclareMathOperator{\Bir}{Bir}
\DeclareMathOperator{\codim}{codim}
\def\UrlBreaks{\do\/\do-\do.}
\title{Nef Abundance: Framework -- K-Trivial, General Dimension}
\date{Temporary draft -- October 11, 2026}
\begin{document}
\maketitle


\section{Theorem and proof}

Fix an integer $d\ge2$. We use the following source inputs,
denoted collectively by $(\mathsf I_d)$.
\begin{itemize}
\setlength{\itemsep}{2pt}
\setlength{\parsep}{0pt}
\item For every projective klt $\Q$-pair $(Y,B)$ of dimension
$j<d$ and every nef $\mathbb{Q}$-Cartier divisor $M$, if $K_Y+B$ is
pseudo-effective and, for some $a\in\Q_{>0}$,
\[
 (K_Y+B+aM)\cdot C>0
\]
for every integral curve through a sufficiently very general point,
then $K_Y+B+cM$ is big for every $c\in\Q_{>0}$.
\item Every projective log canonical $\Q$-pair of dimension at most
$d$ with nef adjoint has semiample adjoint.
\item Every projective $\Q$-factorial generalized klt $\Q$-pair of
dimension at most $d$, with fixed nef $\mathbb{Q}$-Cartier b-divisor
data and pseudo-effective adjoint, admits a terminating choice of
minimal-model program with a nef endpoint. The program extracts no
divisors, preserves $\Q$-factorial generalized klt models, and on a
common resolution the initial adjoint equals the final adjoint plus
an effective divisor exceptional over the final model.
\end{itemize}
The first input is $(P_j)$ of \cite[Section~3.1]{NS}. The last two
are the companion inputs used in \cite[Section~2]{NS}.

\begin{theorem}
Assume $(\mathsf I_d)$. Let $X$ be a projective $\Q$-factorial
terminal variety over $\C$ of dimension $d$, with $K_X\sim0$.
Let $D$ be a nef Cartier divisor with $\nef(X,D)=d$. Then $D$
is big. In particular,
\[
 H^0(X,\OO_X(mD))\ne0
 \qquad\text{for some integer }m>0.
\]
\end{theorem}

\begin{proof}
Suppose for contradiction that $D$ is not big. Since $D$ is
nef, $D^d=0$.

\begin{setup}
\leavevmode
\begin{itemize}
\setlength{\itemsep}{2pt}
\setlength{\parsep}{0pt}
\item Fix an ample Cartier divisor $A$ on $X$.
\item Fix $\epsilon\in\Q_{>0}$ sufficiently small. For $n\ge1$, put
\[
 a_n=((D+A/n)^d)^{1/d},\qquad
 r_n=\max\{1,\lfloor\epsilon/a_n+1/2\rfloor\},\qquad
 P_n=r_n(D+A/n).
\]
Then $P_n$ is ample and $\mathbb{Q}$-Cartier, $r_n\to\infty$, and
$P_n^d\to\epsilon^d$. Indeed, $a_n\to0$ and
$|r_na_n-\epsilon|\le a_n/2$ for all sufficiently large $n$.
The sequence index $n$ is independent of the dimension $d$.
\item Let $p_i:X\times X\to X$ be the projections. Put
\[
 Z=\PP_{X\times X}
       \bigl(p_1^*\OO_X(D)\oplus p_2^*\OO_X(D)\bigr),
 \qquad \pi:Z\to X\times X.
\]
Let $S_1,S_2$ be the two summand sections, and set
\[
 U=\pi^{-1}(X_{\mathrm{reg}}\times X_{\mathrm{reg}})
       \setminus(S_1\cup S_2).
\]
Thus $U$ is smooth and $\dim Z=2d+1$.
\item Set $\xi=c_1(\OO_Z(1))$. For an integer $s>0$, put
\[
 R_{n,s}=(p_1\circ\pi)^*P_n+(p_2\circ\pi)^*P_n+s\xi.
\]
We use the convention
$\pi_*\OO_Z(a)=\Sym^a(p_1^*\OO_X(D)\oplus p_2^*\OO_X(D))$
for $a\ge0$. Whenever $|kR_{n,s}|$ is used, $k>0$ is an integer
with $kP_n$ Cartier.
\item Choose $\epsilon$ sufficiently small that
\[
 (14\epsilon)^d<\tfrac14,\qquad80\epsilon<\tfrac34.
\]
The integer $s$ is chosen after $\epsilon$ and before $n$ tends
to infinity. The multiple $k$ is chosen after $n$.
\end{itemize}
\end{setup}

\medskip
\noindent\textbf{Jet separation.}
Let $V$ be a variety, let $B$ be a Cartier divisor, and let
$z\in V_{\mathrm{reg}}$. For $j\ge0$, $|B|$ separates $j$-jets
at $z$ if
\[
 \forall u\in\OO_{V,z}(B)\quad
 \exists\sigma\in H^0(V,\OO_V(B))\quad
 \sigma_z-u\in\mathfrak m_z^{j+1}\OO_{V,z}(B).
\]

Under the above assumptions, we use the following two lemmas.

\begin{lemma}\label{lem:general-jets}
Assume $(\mathsf I_d)$. Let $X,D$ and $R_{n,s}$ be as in Setup.
There exists an integer $s>0$ such that, for every $n\gg0$,
$|kR_{n,s}|$ separates $(2d+2)k$-jets at a very general point
$z\in U$, for all sufficiently divisible positive integers $k$.
\end{lemma}

This is the large-Seshadri-constant construction of
\cite[Section~5]{NS}, followed by the jet conversion in
\cite[Section~6.7]{NS}.

\begin{lemma}\label{lem:general-nojets}
Assume $(\mathsf I_d)$. Let $X,D$ and $R_{n,s}$ be as in Setup.
For every integer $s>0$ and every $n\gg0$, $|kR_{n,s}|$ does not
separate $(2d+2)k$-jets at a very general point $z\in U$, for any
positive integer $k$ with $kP_n$ Cartier.
\end{lemma}

The terminal application in \cite[Section~6.7]{NS} establishes
the obstruction on a smooth resolution $w:\overline X\to X$.
There $K_{\overline X}\sim E$ for an effective $w$-exceptional
divisor $E$. An ample perturbation of $w^*P_n$ supplies the
complete-intersection flag, and the jets transfer on the
isomorphism locus. For one fixed collection of complex data,
the Frobenius determinant satisfies
\[
 \frac34p^{2d}\le\ord(\sigma)
 \le\bigl(80\epsilon+O(p^{-1})\bigr)p^{2d},
\]
which is impossible for arbitrarily large residual
characteristics $p$.

Fix $s$ supplied by Lemma~\ref{lem:general-jets}, and choose
$n$ sufficiently large for both lemmas. For a sufficiently
divisible $k$, Lemma~\ref{lem:general-jets} gives separation at a
very general point of $U$, whereas
Lemma~\ref{lem:general-nojets} forbids it. This contradiction
proves that $D$ is big.
\end{proof}

\medskip
\noindent\textit{Remark.}
The condition $K_X\sim0$ is imposed at the outset. No adjoint
reduction, index-one cover or terminalization is needed before
Setup, and no irregularity assumption is imposed. Smooth
projective varieties with $K_X\sim0$ are included.
\medskip
\noindent\textit{From nonvanishing to bigness.}
If nonvanishing is already known, bigness also follows from
the log-abundance input in $(\mathsf I_d)$. Choose an effective
$\Q$-divisor $G\sim_{\Q}D$ and a sufficiently small rational
$\tau>0$ such that $(X,\tau G)$ is klt. Then
\[
 K_X+\tau G\sim_{\Q}\tau D\quad\text{is nef}.
\]
Log abundance makes $D$ semiample. A morphism defined by a
free multiple of $D$ has zero-dimensional general fibre:
otherwise a curve in a fibre through a very general point
would have $D$-degree zero, contrary to $\nef(X,D)=d$.
Thus the morphism is generically finite onto its image, and
$D$ is big.

\medskip
\noindent\textbf{Draft status.}
This is a conditional framework following \cite[Sections~5--6]{NS}.
The source inputs $(\mathsf I_d)$ and the inherited geometric and
Frobenius arguments have not been independently verified here.
The two jet lemmas record the source proof obligations; their
complete proofs are not supplied by this framework.

\begin{thebibliography}{99}
\bibitem{NS}
\emph{Numerical Semiampleness of Nef Adjoint Divisors},
OpenAI mathematics collection, October 3, 2026,
Sections~3--6.
\href{https://github.com/openai/math/blob/main/preprints/Numerical-Semiampleness-of-Nef-Adjoint-Divisors-October-3-2026/paper.pdf}{Source manuscript}.
\end{thebibliography}

\end{document}
